\documentclass[11pt,a4paper]{article}
\usepackage[utf8]{inputenc}
\usepackage[T1]{fontenc}
\usepackage{amsmath,amssymb,amsthm}
\usepackage{listings}
\usepackage{geometry}
\usepackage{hyperref}
\geometry{margin=2.5cm}

\newtheorem{lemma}{Lemma}

\lstset{
  language=C++,
  basicstyle=\ttfamily\small,
  keywordstyle=\bfseries,
  breaklines=true,
  numbers=left,
  numberstyle=\tiny,
  frame=single,
  tabsize=2
}

\title{IOI 2019 -- Sequence (Line)}
\author{Solution}
\date{}

\begin{document}
\maketitle

\section{Problem Summary}
Given an array of $n$ integers, find the minimum number of elements to change so the result is an arithmetic sequence $a + bi$ for positions $i = 0, 1, \ldots, n-1$.

\section{Solution}

\subsection{Key Observation}
If we keep elements at positions $i$ and $j$ ($i < j$), the common difference is $d = (\text{arr}[j] - \text{arr}[i]) / (j - i)$, which must be an integer. The starting value is $a = \text{arr}[i] - d \cdot i$.

\begin{lemma}
If the optimal solution keeps $\ge 3$ elements, then at least two of the first three positions ($0, 1, 2$) are kept.
\end{lemma}
\begin{proof}
If none of positions $0, 1, 2$ is kept, at most $n - 3$ elements are kept, which can only be optimal when $n \le 5$. If exactly one is kept, any single element can be paired with any other to define a candidate, but we need at least two kept elements to fix the arithmetic sequence. The lemma follows from a case analysis: among any three consecutive kept positions, at least one pair involves position $0$, $1$, or $2$.
\end{proof}

\subsection{Algorithm}
\begin{enumerate}
  \item For each pair $(i, j)$ where $i \in \{0, 1, 2\}$ and $j > i$: compute $d$ and $a$, then count matches in $O(n)$.
  \item This gives $O(n)$ candidate sequences, each checked in $O(n)$ time.
  \item Handle edge cases: $n \le 2$ always has answer $0$.
\end{enumerate}

\section{Implementation}

\begin{lstlisting}
#include <bits/stdc++.h>
using namespace std;

int minimum_changes(vector<int>& arr) {
    int n = arr.size();
    if (n <= 2) return 0;

    int best = 2;  // can always keep at least 2 elements

    auto try_seq = [&](long long a, long long d) {
        int cnt = 0;
        for (int i = 0; i < n; i++)
            if ((long long)arr[i] == a + d * i) cnt++;
        best = max(best, cnt);
    };

    for (int i = 0; i < min(n, 3); i++) {
        for (int j = i + 1; j < n; j++) {
            long long diff = (long long)arr[j] - arr[i];
            long long gap = j - i;
            if (diff % gap != 0) continue;
            long long d = diff / gap;
            long long a = (long long)arr[i] - d * i;
            try_seq(a, d);
        }
    }

    return n - best;
}

int main() {
    int n;
    scanf("%d", &n);
    vector<int> arr(n);
    for (int i = 0; i < n; i++) scanf("%d", &arr[i]);
    printf("%d\n", minimum_changes(arr));
    return 0;
}
\end{lstlisting}

\section{Complexity Analysis}
\begin{itemize}
  \item \textbf{Time}: $O(n^2)$ -- $O(n)$ candidate pairs, each checked in $O(n)$.
  \item \textbf{Space}: $O(n)$.
\end{itemize}

\end{document}
